\chapter{Discussion}
\label{ch:discussion}

\section{Evaluation of the Hypotheses}
\label{sec:disc-hyp}

\paragraph{H1 (reduced-order optimization)} is supported within the model. Whenever the model contains lift and it is not
suppressed, all 101 feasible optima are lift-propelled. H1 is tested not only across models, where the lift-propelled gaits
need one half to one third of the power of the best drag-only gaits, but also within one model: drag-based strokes forced by
a lift-share limit need four to ten times the power and do not reach \SI{0.30}{\metre\per\second}, although this second factor
rests on only 3 feasible strokes. The statement is conditional on an uncalibrated lift model, as discussed below.

\paragraph{H2 (drag versus lift across speed)} is supported. As expected, the optimized paddle
is much stronger at low speed and its thrust falls steeply with speed, to about zero at \SI{0.30}{\metre\per\second}; the
fluke keeps a larger fraction of its thrust and becomes stronger above \SIrange{0.23}{0.24}{\metre\per\second}. The cause of
the paddle's speed limit is identified: the recovery drag, which grows more than sevenfold between rest and
\SI{0.30}{\metre\per\second}, and the momentum needed to open the fan in the oncoming flow, while the power-stroke thrust falls
by less than half. The efficiency crossover lies slightly below the thrust crossover, at about \SI{0.22}{\metre\per\second}; below it, the optimized paddle is both stronger and more efficient than the fluke, reaching $\etaF=0.35$--0.37 at \SI{0.15}{\metre\per\second}. This differs from the biological picture, in which lift-based swimmers are more efficient than paddlers at all speeds \cite{Fish1996}; for a small robot leg with a well-folding paddle it holds only above the crossover. The second part of H2 is confirmed: scheduling the fluke
pitch so that the angle of attack stays near \SI{20}{\degree} raises the efficiency at \SI{0.30}{\metre\per\second} from 0.24 to 0.40.

The paddle's efficiency exceeds the maximum of about 0.33 that Fish reported for paddling animals \cite{Fish1984,Fish1996}. The two
values are not directly comparable: the animal values are whole-animal estimates, while $\etaF$ here is the hydrodynamic Froude
efficiency of an isolated propulsor whose folded fan produces only about one ninth of the force of the open fan ($a=0.11$, against an area ratio of 0.445 for the muskrat). It also
stays below the quasi-steady bound $\U/v_p\approx0.6$ at \SI{0.15}{\metre\per\second}, with the mean power-stroke speed
$v_p\approx\SI{0.25}{\metre\per\second}$ of the fan centroid.

\paragraph{H3 (whole robot)} is largely supported. With the best-estimate resistance the fluke reaches \SI{0.30}{\metre\per\second}
and the paddle \SIrange{0.26}{0.27}{\metre\per\second}, so the speeds are similar but not equal, and the speed-scheduled fluke needs
43\,\% less power than the paddle (\SI{42}{} versus \SI{73}{\milli\watt}). A first hardware test, however, is 1.3 to 1.6 times slower than predicted (\cref{sec:res-hw}). The predicted speeds are therefore upper bounds for the present robot, and the comparison
of the propulsors rests on the CFD.

\section{Why the Reduced-Order Model and the CFD Differ at Low Speed}
\label{sec:disc-models}

The reduced-order model and the CFD agree that drag-based paddling has a speed limit and that lift-based flapping wins at cruising
speed. They disagree at low speed: MuJoCo finds lift-propelled gaits more economical already at
\SI{0.10}{\metre\per\second}, by a factor of about two to four (up to ten at \SI{0.20}{\metre\per\second}), whereas the CFD finds the optimized paddle more efficient below about
\SI{0.22}{\metre\per\second}. Four differences explain this.
\begin{enumerate}
  \item \emph{The lift model is optimistic.} The flat-plate law $c_l\sin2\alpha$ produces lift without the induced drag and
        tip losses of a finite wing, and the flipper of the toy model has an aspect ratio of only 1.4. The validation case suggests that even at an aspect ratio of 6 the finite span is a likely cause of an efficiency deficit of up to 18\,\%, with tip losses estimated at 10--20\,\%. The model also has no dynamic stall,
        no delay in the build-up of circulation and no interaction of the flipper with its own wake.
  \item \emph{The drag model is pessimistic for a starting paddle.} A constant $C_D=2.2$ misses the starting vortex and added-mass
        effects that raise the effective force coefficient of the paddle's power stroke to 2--3 times its quasi-steady value in
        the CFD (\cref{sec:res-drag}).
  \item \emph{The drag-type strokes differ.} The MuJoCo flipper can only feather by rotating; it has no folding and no controlled
        fold timing, and it is limited to the smooth harmonic motion \eqref{eq:gait}. Because a flipper that crosses the water at an
        angle produces lift, a drag-based stroke in the lift model is a stroke from which lift has been suppressed, and most
        paddling strokes of the drag-only model become lift-propelled once lift is added (\cref{sec:mj-liftshare}). The CFD
        paddle, in contrast, is a dedicated drag device: it moves broadside to its path, folds to 12\,mm at optimized phases and
        uses a trapezoidal velocity law, each feature being worth up to about 30\,\% in thrust. MuJoCo thus compares lift-based strokes with
        lift-suppressed flipper strokes, the CFD a fluke with an optimized folding paddle.
  \item \emph{The optimization effort is asymmetric in both directions.} In MuJoCo the lift gaits were optimized for power at a
        required speed, while the drag-only front has not converged and the constrained search for drag-based strokes in the
        lift model found a feasible stroke in only 3 of 30 runs. In the CFD, the paddle was optimized in several rounds,
        while the fluke motion L0 was taken from another model; only its pitch amplitude was later scheduled with speed. The
        areas also differ (\SI{18.5}{} versus \SI{11.2}{\square\centi\metre}).
\end{enumerate}
The reduced-order model is therefore best read as an argument about \emph{mechanism}, which propulsion principle an
optimizer exploits when it can, and the CFD as the argument about \emph{magnitude} and about the speed at which one
principle overtakes the other. Where the two disagree, the CFD, which resolves the flow and is validated, is taken as
authoritative. A calibration of the reduced-order coefficients, especially $c_l$ with an aspect-ratio correction, against
the CFD would reconcile the two and is the natural next step.

\section{Implications for BODY2}
\label{sec:disc-design}
\label{sec:disc-limits-prop}

\paragraph{Recommendation.} For crossing open water at the surface, the recommended propulsor for BODY2 is the fluke with a pitch
that decreases with speed. For the best-estimate resistance it is slightly faster than the optimized paddle V3trap, needs 43\,\% less
hydrodynamic power, and requires neither a controlled fold nor a trapezoidal velocity law. Compared with the robot's present
configuration, it roughly triples the predicted cruising speed, and in an untimed comparison on the hardware the flukes swam visibly faster than the paddles
(\cref{sec:plat-hw}). The
paddle remains the better choice where acceleration, station keeping at low speed or walking matter; it then needs a fold narrowed to
about \SI{12}{\milli\metre} that can be commanded within \SIrange{45}{60}{\milli\second} at a specific phase, which the current cable
fold cannot do. This division of labour, rowing for starting and low speed, flapping for cruising, is the one found in animals
\cite{Walker2000,Fish1996}; a leg that could change from paddle to fluke, like the morphing turtle limbs of Baines et
al.\ \cite{Baines2022}, would cover both regimes.

\paragraph{Fluke.} The fluke's thrust at rest is limited because it is stalled: without forward speed, a flapping foil acts
as a small paddle. At speed, its performance depends on keeping the angle of attack near the efficient range, which a fixed
pitch schedule does not do. The schedule $\theta_0=\gamma_\text{max}/2$ is a single-parameter rule that achieves this and also
removes the large mean vertical force that a mismatched schedule produces. On the hardware the pitch is passive: it follows from
the balance between the hydrodynamic moment about the pin, which turns the fluke towards the inflow, and the leaf spring. A passive
fluke therefore adapts to speed, but not in the ratio required. The hydrodynamic moment grows with the square of the relative speed
while the spring moment does not, so the ratio $\theta/\gamma$ rises with speed. A linear spring can match
$\theta_0=\gamma_\text{max}/2$ at one design speed only; below it the fluke pitches too little and its angle of attack is too large,
above it the fluke aligns too far with the inflow and its angle of attack falls below the efficient range. Matching a range of
speeds requires a progressive spring, an adjustable preload or an active pitch actuator. The hardware test, which reached only
roughly a quarter to a half of the CFD thrust (\cref{sec:res-hw}), is consistent with a passive pitch that is not tuned to the efficient angle of attack, although
other causes cannot be excluded.

\paragraph{Equal area.} The fluke used here has 61\,\% of the paddle's area. At rest both produce the same thrust per area
(\SIrange{3.9}{4.5}{} versus \SI{4.1}{\milli\newton\per\square\centi\metre}), and at \SI{0.223}{\metre\per\second} the fluke already
produces more (\SI{1.9}{} versus \SIrange{1.3}{1.4}{\milli\newton\per\square\centi\metre}). If the thrust and power coefficients stayed
the same, a fluke of the paddle's area, \SI{18.5}{\square\centi\metre}, moved with the same kinematics would produce about
\SI{75}{\milli\newton} at rest, as much as the paddle, and more than the paddle from about \SI{0.15}{\metre\per\second} on. Its power would
grow in the same proportion, to about \SI{30}{\milli\watt} at rest, so the efficiency crossover would not move. A lunate fluke of this area
(span \SI{86}{\milli\metre}, $\mathit{AR}\approx4$) has been designed for the robot; its larger aspect ratio should reduce tip losses, but its
larger hinge moment, about 1.6 times that of A2 at the same pitch schedule, must be carried by the pitch spring or actuator. This is an
estimate, not a simulation result.

\paragraph{Actuators.} The hydrodynamic moments of all motions are small compared with the servos: the largest, \SI{36}{\milli\newton\metre}
about the hip for the trapezoidal paddle stroke, is 6--8\,\% of the stall torque of the PTK\,7465W (\SIrange{0.44}{0.57}{\newton\metre}).
The limits that matter are joint speed and the ability to switch the fold (\cref{tab:actuators}). The prescribed pitch of L0 needs
\SI{712}{\degree\per\second}, more than the no-load speed at \SI{6}{\volt} (\SI{650}{\degree\per\second}) and close to that at
\SI{8.4}{\volt} (\SI{860}{\degree\per\second}), whereas the speed-scheduled pitch at \SI{0.30}{\metre\per\second} needs only \SI{366}{\degree\per\second}.

\section{Limitations}
\label{sec:disc-limits}

\begin{itemize}
  \item \emph{Single leg.} The CFD simulates one propulsor without the linkage, the hull or the other legs. Interaction between
        legs can raise the thrust of four paddling legs above four times the single-leg value \cite{Wang2025} and may differ between
        paddle and fluke. The self-propulsion balance neglects it.
  \item \emph{Free surface.} All propulsor CFD is single-phase with a rigid lid. On the hardware and in Flare Live the fluke breaks the
        surface by about \SI{7}{\milli\metre} at the top of its stroke, and the paddle's recovery path is close to it; ventilation
        and waves are not modelled.
  \item \emph{Prescribed pitch and fold.} The fluke pitch is prescribed rigidly, while the hardware fluke is passive. The fan's
        folding is modelled as an instantaneous switch between two separately simulated states. The change of added mass at the
        switches is corrected analytically within bounds, with an added mass identified from the CFD; the effect of the open run's
        own recovery wake on its power stroke is not corrected and makes the paddle results slightly optimistic.
  \item \emph{CFD accuracy.} The flow is laminar, which is one of the causes of the 18\,\% efficiency deficit in the validation at
        $\mathit{Re}=4\times10^4$; at the lower Reynolds numbers of the fluke the effect is expected to be smaller. The fluke thrust carries
        about $\pm12$\,\% discretization uncertainty. The paddle mesh is coarser and was not studied separately; the conclusions were argued in \cref{sec:ver-drag} to be insensitive to an error of that size.
  \item \emph{Power.} The power is the hydrodynamic input power of the propulsor. Servo losses, the inertia of the leg and the drag of
        the linkage are not included, so the power ratios between propulsors are not the ratios of battery power.
  \item \emph{Unequal propulsors.} The paddle stroke was optimized more thoroughly than the fluke motion, so the fluke results are
        conservative; a systematic optimization of heave amplitude, frequency and pitch would likely move the crossover to lower speed.
        The propulsors also run at different frequencies (\SI{1.39}{} and \SI{2}{\hertz}) and have different areas; thrust per area and
        per power are reported, but no comparison at equal area and frequency was made.
  \item \emph{Resistance.} The hull resistance is known only within bounds, and the cruising speeds are correspondingly
        uncertain. The linkage term of $D_\text{mid}$ is an engineering estimate, not a measurement, and the hull resistance was computed
        for the pointed end first, whereas the hardware swims with its blunt end first.
  \item \emph{Reduced-order model.} The MuJoCo model is a stand-in of similar size, not BODY2, and its coefficients are not
        calibrated. Its servo model (\SI{3}{\newton\metre}, \SI{400}{\degree\per\second}) is stronger and slower than the PTK\,7465W
        (\SIrange{0.44}{0.57}{\newton\metre}, \SIrange{650}{860}{\degree\per\second}). Its searches are not fully converged, and the
        within-model comparison rests on 3 feasible drag-based strokes from 30 searches with two seeds, so its power factor may
        be overestimated.
  \item \emph{Hardware.} The hardware test consists of one hand-timed run per frequency from rest, with a reversed hull, a passive
        pitch and an unmeasured stroke amplitude. It shows that the self-propulsion estimate is an upper bound but does not identify
        the causes of the difference.
\end{itemize}
